\bigskip
Das \textit{Tagebuch} führt nunmehr in den Juli des Jahres 1798, nämlich zu
den Notizen 91a., 91b., 92. — Die in 91a. angegebene Reihe findet sich
auch in der Leistenotiz Werke X, 1, S. 169, Art. 5, wo sie zur Berechnung
von \(\frac{\varpi}{\pi}\) auf 15 Dezimalstellen benutzt wird. In 91b. tritt ein wesentlich neuer
Gedanke auf (vergl. Werke X, 1, S. 168, Art. 4), nämlich die Entwicklung von
\(\operatorname{sin\,lemn}a\) in eine Reihe, die nach den Sinus der ungeraden Vielfachen von \(a\)
fortschreitet, und eine Entwicklung des \(\operatorname{arc\,sin\,lemn}\sin\varphi\) nach den Sinus der
geraden Vielfachen des Winkels \(\varphi\) mit in Graden ausgedrückten Zahlkoeffizienten.
Dieser Gedanke bildet die Vorbereitung zu den bedeutungsvollen
Entdeckungen, die die Tagebuchnotiz 92. anzeigt, wo \textsc{Gauss} sagt, er habe
»von der Lemniskate alle Erwartungen übertreffende Feinheiten gefunden, und
zwar durch Methoden, die uns ein wahrhaft neues Feld eröffnen«. Wir finden
diese Resultate in der gleichfalls im Juli 1798 begonnenen \textls[120]{Scheda Aa}, dem
ersten der von \textsc{Gauss} als \textls[120]{Schedae} bezeichneten Notizheftchen\footnote{Aus dieser sind Werke III abgedruckt der Aufsatz \textit{De curva lemniscata} S. 413ff., u. z. Artikel 1,
2, 3 (S. 3—4 der Scheda), [4] (S. 6 der Scheda), [6] (Zeile 2—6, S. 7 der Scheda, Zeile 7, 8, Seite 8 der
Scheda), [7] (bis Ende der S. 417, S. 14, 15 der Scheda, die beiden ersten Gleichungen auf S. 418, S. 21
und 22 der Scheda), [8] (die erste Gleichung S. 27, die vier folgenden und Zeile 2—4 v. u. S. 28 der Scheda).
Die Angabe \textsc{Scherings} (Werke III, S. 494), dass dieser Art. [8] aus der Scheda Ac stammt, beruht auf
einem Versehen.}. Diese beginnt
mit einer augenscheinlich das früher Gefundene zusammenfassenden Redaktion
der Untersuchungen über lemniskatische Funktionen, die aber bald
(Seite 4) abbricht. Die folgenden Einträge (S. 6—8 der Scheda) zeigen wieder
ganz den Charakter vorläufiger Notizen und geben die in der Tagebuchauf-\par
\clearpage\SourceLeaf{123}\ArithmeticHead{38}{L}{L. SCHLESINGER, ÜBER GAUSS’ FUNKTIONENTHEORETISCHE ARBEITEN.}
\noindent zeichnung 92. bezeichneten Entdeckungen wieder, was in den Bemerkungen
zu dieser Aufzeichnung (Werke X, 1, S. 535—536) im einzelnen belegt ist.
Hier geben wir nur eine Übersicht über die erzielten Fortschritte. Es werden
für die Zähler und Nenner der Funktionen sinus und cosinus lemniscaticus
an der Stelle der früher benutzten \(M,N,m,n\) die den \textsc{Jacobi}schen Thetafunktionen
entsprechenden Funktionen \(P,Q,p,q\) eingeführt, die sich von
jenen durch Exponentialfaktoren unterscheiden; z. B. ist (vergl. Werke X, 1,
S. 168, Art. 3 und S. 170, Art. 6)
\[
P(x)=M(x)\,.\,e^{-\frac12\frac{x^2}{\varpi^2}\pi}.
\]
\noindent Sie erscheinen in der schon im \textit{Tagebuch} 91b. eingeführten Variabeln \(\sin\varphi\)
vermöge der Kenntnis der Null- und Unendlichkeitsstellen von sin lemn als
\textls[120]{einfach} unendliche Produkte, und diese Darstellung liefert sofort den Beweis
für die im \textit{Tagebuch} 63. am 29. März 1797 angezeigte Gleichung \(N(\varpi)=e^{\frac12\pi}\)
(vergl. oben). Ferner ergeben sich die Koeffizienten der Entwicklung von
\(\operatorname{sin\,lemn}\varphi\) nach den Sinus der ungeraden Vielfachen von \(\varphi\) in allgemeiner Form:
\begingroup\makeatletter\tagsleft@true\let\veqno\leqno\makeatother
\begin{equation*}\tag*{(1)}
\operatorname{sin\,lemn}\varphi=\frac{\pi}{\varpi}\frac4{e^{\frac12\pi}+e^{-\frac12\pi}}\sin\varphi
-\frac{\pi}{\varpi}\frac4{e^{\frac32\pi}+e^{-\frac32\pi}}\sin3\varphi+\ldots,
\end{equation*}
\endgroup
\noindent während sie im \textit{Tagebuch} 91b. nur numerisch gegeben waren.

Um nun den letzten Schritt zu vollziehen, der in der Darstellung von
Zähler und Nenner selbst in der Form von trigonometrischen Reihen besteht,
stellt \textsc{Gauss} zunächst sehr mühsame und scharfe numerische Rechnungen für
\(\frac{\pi}{\varpi},e^\pi,e^{-\pi},e^{\frac12\pi},e^{-\frac12\pi}\) an, die zum Teil Werke III, S. 431 Zeile 7 v. u. bis
S. 432 abgedruckt sind\footnote{\textsc{Gauss} hat bei diesen Rechnungen Methoden der additiven Zahlentheorie zur Anwendung
gebracht, namentlich die Darstellung einer Zahl durch die Summe von zwei Quadraten, um ihren
Primzahlcharakter zu entscheiden. Solche Rechnungen finden sich vielfach auch schon im \textit{Leiste}, so S. 48:
»discerpendi 283009 in bina quadrata \(4225+528^2\), Also Primzahl«, ferner S. 96, 103 usw. — Die Werke
III, S. 426—431 abgedruckten, auf wesentlich feineren Methoden beruhenden Rechnungen für dieselben
Grössen entstammen einer späteren Zeit, da darin die \textit{Disqu. Arithm.} zitiert werden (allerdings ohne Angabe
der Artikelnummer, so dass es möglich wäre, dass das genannte Werk damals noch nicht ganz fertig
gedruckt war), wahrscheinlich dem Jahre 1800 im Zusammenhang mit der Tagebuchnotiz 112.; vergl. weiter
unten Abschn. V c).}, und verfährt dann in folgender Weise weiter:

Es werden für diese Zähler und Nenner \(P\varphi,Q\varphi,p\varphi,q\varphi\) die Logarithmen
genommen, die sich aus der Darstellung der Funktionen selbst in Form von\par
\clearpage\SourceLeaf{124}\ArithmeticHead{39}{R}{III b. LEMNISKATISCHE FUNKTIONEN 1796—1799.}
\noindent einfach unendlichen Produkten mit Benutzung der Formel\footnote{Scheda Aa, S. 14, Werke III, S. 417; vergl. \textsc{Euler}, \textit{calc. integralis} I, § 295, Opera, ser. I, vol.
11, S. 179.}
\[
\log(1+\mu\cos\varphi)=2\sum_{n=1}^{\infty}(-1)^{n+1}
\left\{\frac{\mu}{1+\sqrt{1-\mu^2}}\right\}^n\frac{\cos n\varphi}{n}
\]
direkt als Reihen ergeben, die nach den Kosinus der geraden Vielfachen von
\(\varphi\) fortschreiten. Die Koeffizienten dieser Reihen werden dann (Scheda Aa,
S. 21, 22, Werke III, S. 418 oben) numerisch angegeben. Nun knüpft \textsc{Gauss}
(Scheda Aa, S. 23) an die auch im \textit{Leiste} S. 78 (Werke X, 1, S. 167, Art. 2)
aufgezeichneten Differentialbeziehungen der \(P\varphi\) usw. an und findet die Darstellung
von \(P\varphi\) selbst in Form einer Reihe, die nach Sinus der Vielfachen
von \(\varphi\) fortschreitet, zunächst mit numerischen Koeffizienten\footnote{Die betreffende Stelle der Scheda Aa fehlt in Werke III, sie ist Werke X, 1, S. 537 wiedergegeben.}.
Es tritt dann
— offenbar als Ergebnis einer numerischen Induktion — auf S. 27 der Scheda
Aa die Formel
\begingroup\makeatletter\tagsleft@true\let\veqno\leqno\makeatother
\begin{equation*}\tag*{(2)}
\operatorname{sin\,lemn}\varphi=
\mathord{\vcenter{\hbox{\Large$\surd$}}}\frac4{e^{\frac{\pi}{2}}}\, .\,
\frac{\sin\varphi-e^{-2\pi}\sin3\varphi+e^{-6\pi}\sin5\varphi-\ldots}
{1+2e^{-\pi}\cos2\varphi+2e^{-4\pi}\cos4\varphi+\ldots}
\end{equation*}
\endgroup
\noindent auf, der sich dann auf S. 28 eben jene trigonometrischen Reihen für \(P,Q\)
anschliessen:
\begingroup\makeatletter\tagsleft@true\let\veqno\leqno\makeatother
\begin{equation*}\tag*{(3)}
\left\{\begin{aligned}
P\varphi&=2^{\frac34}\sqrt{\frac{\pi}{\varpi}}\left\{
\frac{\sin\varphi}{e^{\frac{\pi}{4}}}-\frac{\sin3\varphi}{e^{\frac{9\pi}{4}}}
+\frac{\sin5\varphi}{e^{\frac{25\pi}{4}}}-\ldots\right\}\\[1ex]
Q\varphi&=\frac1{2^{\frac14}}\sqrt{\frac{\pi}{\varpi}}\left\{
1+\frac{2\cos2\varphi}{e^\pi}+\frac{2\cos4\varphi}{e^{4\pi}}+\ldots\right\}\hbox{\footnotemark}
\end{aligned}\right.
\end{equation*}
\endgroup
\footnotetext{Die Gleichungen (1) bis (3) sind hier getreu nach der Handschrift abgedruckt; \textsc{Gauss} denkt \(\varphi\) auf
beiden Seiten im Gradmass gegeben. Will man, wie \textsc{Schering} beim Abdruck in Werke III getan hat, auf
Bogenmass reduzieren, so muss man linker Hand, d. h. für die Lemniskate, \(180^\circ=\varpi\) und rechter Hand,
d. h. für den Kreis, \(180^\circ=\pi\) nehmen. Wenn also \(\varphi^\circ/180^\circ=\psi\) gesetzt wird, so erhält man in der Tat,
wie Werke III gedruckt ist, links \(\psi\varpi\), rechts \(\psi\pi\).}
\noindent und die daraus folgenden:
\begingroup\makeatletter\tagsleft@true\let\veqno\leqno\makeatother
\begin{equation*}\tag*{(4)}
\left\{\begin{aligned}
1-2e^{-\pi}+2e^{-4\pi}-2e^{-9\pi}+\ldots&=\sqrt{\frac{\varpi}{\pi}}\\[1ex]
e^{-\frac{\pi}{4}}+e^{-\frac{9\pi}{4}}+e^{-\frac{25\pi}{4}}+\ldots&=\frac12\sqrt{\frac{\varpi}{\pi}}.
\end{aligned}\right.
\end{equation*}
\endgroup
\clearpage\SourceLeaf{125}\ArithmeticHead{40}{L}{L. SCHLESINGER, ÜBER GAUSS’ FUNKTIONENTHEORETISCHE ARBEITEN.}
Wir sehen hier zunächst in (3) die Grösse \(\frac{\pi}{\varpi}\) auftreten, die \textsc{Gauss} (vergl.
oben S. 32) so vielfach numerisch berechnet hat. Ferner könnte die Darstellung
(4) von \(\sqrt{\frac{\varpi}{\pi}}\) durch Potenzreihen von \(e^{-\pi}\) bezw. \(e^{-\frac14\pi}\), deren Exponenten
eine arithmetische Progression zweiter Ordnung bilden (die Quadratzahlen
sind), eine erste Ahnung des Zusammenhanges suggeriert haben, der
zwischen der Grösse \(\frac{\varpi}{\pi}\) und dem agM. besteht.

Man geht wohl nicht fehl, wenn man die zuletzt geschilderten Entdeckungen,
die die Theorie der lemniskatischen Funktionen zur höchsten
Stufe der Entwicklung bringen, mit der Tagebuchnotiz 95. vom Oktober
1798 in Verbindung setzt, der ersten die \textsc{Gauss} nach dem Abgang von
der Universität (siehe weiter unten d)) wieder in Braunschweig geschrieben
hat, wo es heisst: »Ein neues Feld der Analysis hat sich uns eröffnet, nämlich
die Untersuchung der Funktionen etc.« Das was sich unausgesprochen
hinter dem etc. verbirgt, wäre eben das Bewusstsein, dass die bei den lemniskatischen
Funktionen angewandten Methoden und die mit ihrer Hilfe zur
Erscheinung gebrachten Eigenschaften, nicht nur etwas durchaus neuartiges
darstellen, wie man es bei den elementaren Transzendenten nicht kennt, sondern
dass die Tragweite dieser Methoden auch nicht auf die so spezielle
lemniskatische Funktion beschränkt sein könne, dass diese vielmehr nur das
erste Glied einer ausgedehnten Klasse von neuartigen Funktionen bildet, die
ein neues Feld der Analysis ausmachen, und zu denen sich \textsc{Gauss} den Zugang
erschlossen hat. In der Überzeugung, sich ein solches Gebiet erobert zu haben,
dürfte ihn der vermutete Zusammenhang zwischen der lemniskatischen Funktion
und dem agM. bestärkt haben, wie daraus hervorgeht, dass die Tagebuchnotiz
98. vom 30. Mai 1799, die diesen Zusammenhang zunächst auf
Grund einer numerischen Induktion ausspricht, mit den zuversichtlichen Worten
endet: »wenn dies bewiesen sein wird, so ist damit \textls[120]{sicher} ein neues Feld
der Analysis erschlossen«. Ehe wir nun dazu übergehen, die Entwicklung
der Theorie des agM. bis zu diesem Zeitpunkte zu schildern, mögen noch
einige allgemeine Bemerkungen über das bereits Dargelegte angefügt werden.

Auf die Entwicklung des \textsc{Gauss}schen Gedankenganges wirft eine Bemerkung
einiges Licht, die \textsc{Gauss} 1828 in einem vom 30. März datierten\par
\clearpage\SourceLeaf{126}\ArithmeticHead{41}{R}{III c. DAS AGM. BIS ZUM MAI 1799.}
\HistoricalFoot{X 2 Abh.2.}{6}
\noindent Briefe an \textsc{Bessel} macht. Er schreibt\footnote{\textit{Gauss-Bessel Briefwechsel}, S. 477, vergl. auch den Werke III, S. 495 zitierten Brief von \textsc{Gauss}
an \textsc{Crelle}.}: »Zur Ausarbeitung der seit vielen
Jahren (1798) angestellten Untersuchungen über die transzendenten Funktionen
werde ich vorerst wohl noch nicht kommen können, da erst noch mit manchen
andern Dingen aufgeräumt werden muss. Herr \textsc{Abel} ist mir, wie ich sehe,
jetzt zuvorgekommen und überhebt mich in Beziehung auf etwa ein Drittel
dieser Sachen der Mühe, zumal, da er alle Entwicklungen mit Eleganz und
Konzision gemacht hat. Er hat gerade denselben Weg genommen, welchen
ich 1798 einschlug, daher die grosse Übereinstimmung der Resultate nicht zu
verwundern ist. .\ .\ . Jeder Missdeutung zuvorzukommen bemerke ich jedoch,
dass ich mich nicht erinnere, von diesen Sachen irgend jemand etwas mitgeteilt
zu haben\footnote{Auf den Einfluss, den \textsc{Gauss} (unmittelbar und mittelbar) auf \textsc{Abel} ausgeübt hat, wird unten
im Abschn. IX a) näher einzugehen sein.}«.

In der Tat ist \textsc{Gauss} ebenso wie \textsc{Abel} zunächst vom Additionstheorem
ausgegangen, hat mit Hilfe desselben die Periodizität bewiesen, dann die Nullstellen
und Unendlichkeitsstellen der Funktionen sin lemn und cos lemn bestimmt
und aus diesen »nach bekannten Regeln« (vergl. Fussnote\textsuperscript{50a)}) die Darstellungen
der Zähler und Nenner in der Form von unendlichen Produkten
gewonnen.

Mit der Entwicklung der Zähler und Nenner in trigonometrische Reihen,
wie überhaupt mit der Einführung dieser als selbständiger Transzendenten, ist
\textsc{Gauss} allerdings schon 1798 erheblich über \textsc{Abels} \textit{Recherches} hinausgegangen,
— wenigstens im lemniskatischen Falle. — Diese Einschränkung muss gemacht
werden, da keinerlei Anhaltspunkt dafür vorliegt, dass \textsc{Gauss} schon
1798 die \textls[120]{allgemeinen} elliptischen Funktionen behandelt hat, und man wird
die Bemerkung von \textsc{Gauss} über seine und \textsc{Abels} Methoden so aufzufassen
haben, dass die 1798 für den Fall der lemniskatischen Funktionen zur Anwendung
gebrachten Methoden gemeint seien, die sich ja dann nachträglich
auch im allgemeinen Falle bewährt haben.
